% GENERATED FILE. Edit canonical lessons, then run npm run generate:books.
% canonical-source-hash: fnv1a64:829225e6657a6e93
% canonical-lessons: PA-C264-par, PA-C264-langhna, PA-C264-phatak, PA-C264-khambha, PA-C264-kona

\chapter{Across the Line}
\label{ch:pa-a2-across-the-line}
\hlchaptermodality{264}

\begin{chapteropening}
\textbf{By the end of this chapter:} \emph{I can say how to get across a road or a line.}

Complete the last lesson of chapter 264: I can say how to get across a road or a line.
\end{chapteropening}

\section[\texorpdfstring{\pa{ਪਾਰ} (pār)}{ਪਾਰ (pār)}]{\pa{ਪਾਰ} (pār) — across; on the far side}

\label{lesson:PA-C264-par}



\begin{quote}
Before the new one: say the Punjabi for right-hand, then the Punjabi for left-hand.
\end{quote}

\subsection*{You'll want to know: \pa{ਪਾਰ}}
\textbf{\pa{ਪਾਰ}} — \emph{pār} — \textquotedblleft{}across; on the far side\textquotedblright{}.

\begin{quote}
\pa{ਸਟੇਸ਼ਨ} \pa{ਪੁਲ} \pa{ਦੇ} \pa{ਪਾਰ} \pa{ਹੈ।}
\end{quote}

\begin{quote}
\emph{saṭeshan pul de pār hai.}
\end{quote}

\begin{quote}
The station is across the bridge.
\end{quote}

\begin{cousinweb}
\textbf{\pa{ਪਾਰ}} (\emph{pār}) is inherited from Sanskrit \emph{pāra}, \textquotedblleft{}further bank, furthest end\textquotedblright{}; CDIAL lists Punjabi \emph{pār}, \textquotedblleft{}beyond\textquotedblright{}, under it (CDIAL 8100).
\end{cousinweb}

\begin{grammarlens}[title={What you've built}]
One more word for this chapter.
\end{grammarlens}

\subsection*{Your turn}
\begin{itemize}
\raggedright
  \item \emph{Say it:} \emph{pār}
  \item \emph{Say it:} \emph{pār}, once more
  \item \emph{Say it:} \emph{saṭeshan pul de pār hai.} — The station is across the bridge.
  \item \emph{Recall:} say the Punjabi for right-hand, then the Punjabi for left-hand, then say \emph{pār} again
\end{itemize}

\begin{tcolorbox}[breakable,colback=teal!4,colframe=teal!35!black,title={Before you move on}]
What does \pa{ਪਾਰ} mean? (\textquotedblleft{}Across; on the far side\textquotedblright{}.) Say it once more.
\end{tcolorbox}
\section[\texorpdfstring{\pa{ਲੰਘਣਾ} (laṅghṇā)}{ਲੰਘਣਾ (laṅghṇā)}]{\pa{ਲੰਘਣਾ} (laṅghṇā) — to pass through; to go past}

\label{lesson:PA-C264-langhna}



\begin{quote}
Before the new one: say the Punjabi for left-hand, then the Punjabi for across.
\end{quote}

\subsection*{You'll want to know: \pa{ਲੰਘਣਾ}}
\textbf{\pa{ਲੰਘਣਾ}} — \emph{laṅghṇā} — \textquotedblleft{}to pass through; to go past\textquotedblright{}.

\begin{quote}
\pa{ਮੈਂ} \pa{ਰੋਜ਼} \pa{ਬਾਜ਼ਾਰ} \pa{ਤੋਂ} \pa{ਲੰਘਦਾ} \pa{ਹਾਂ।}
\end{quote}

\begin{quote}
\emph{maiṁ roz bāzār ton laṅghdā hāṁ.}
\end{quote}

\begin{quote}
I pass through the market every day.
\end{quote}

Past (I …): \pa{ਮੈਂ} \pa{ਪੁਲ} \pa{ਤੋਂ} \pa{ਲੰਘਿਆ}, \emph{maiṁ pul ton laṅghiā}.

\begin{cousinweb}
\textbf{\pa{ਲੰਘਣਾ}} (\emph{laṅghṇā}) is inherited from Sanskrit \emph{laṅghayati}, \textquotedblleft{}leaps over\textquotedblright{}; CDIAL lists Punjabi \emph{laṅghṇā}, \textquotedblleft{}to pass\textquotedblright{}, under it (CDIAL 10905).
\end{cousinweb}

\begin{grammarlens}[title={What you've built}]
One more word for this chapter.
\end{grammarlens}

\subsection*{Your turn}
\begin{itemize}
\raggedright
  \item \emph{Say it:} \emph{laṅghṇā}
  \item \emph{Say it:} \emph{laṅghṇā}, once more
  \item \emph{Say it:} \emph{maiṁ roz bāzār ton laṅghdā hāṁ.} — I pass through the market every day.
  \item \emph{Recall:} say the Punjabi for left-hand, then the Punjabi for across, then say \emph{laṅghṇā} again
  \item \emph{Write it:} copy \textbf{\pa{ਲੰਘਣਾ}} once, checking the whole word
\end{itemize}

\begin{tcolorbox}[breakable,colback=teal!4,colframe=teal!35!black,title={Before you move on}]
What does \pa{ਲੰਘਣਾ} mean? (\textquotedblleft{}To pass through; to go past\textquotedblright{}.) Say it once more.
\end{tcolorbox}
\section[\texorpdfstring{\pa{ਫਾਟਕ} (phāṭak)}{ਫਾਟਕ (phāṭak)}]{\pa{ਫਾਟਕ} (phāṭak) — a big gate; a level-crossing gate}

\label{lesson:PA-C264-phatak}



\begin{quote}
Before the new one: say the Punjabi for across, then the Punjabi for to pass through.
\end{quote}

\subsection*{You'll want to know: \pa{ਫਾਟਕ}}
\textbf{\pa{ਫਾਟਕ}} — \emph{phāṭak} — \textquotedblleft{}a big gate; a level-crossing gate\textquotedblright{}.

\begin{quote}
\pa{ਜਦੋਂ} \pa{ਰੇਲ} \pa{ਆਉਂਦੀ} \pa{ਹੈ} \pa{ਤਾਂ} \pa{ਫਾਟਕ} \pa{ਬੰਦ} \pa{ਹੁੰਦਾ} \pa{ਹੈ।}
\end{quote}

\begin{quote}
\emph{jadoṁ rel āundī hai tāṁ phāṭak band hundā hai.}
\end{quote}

\begin{quote}
When the train comes, the crossing gate is closed.
\end{quote}

A masculine noun.

\begin{cousinweb}
\textbf{\pa{ਫਾਟਕ}} (\emph{phāṭak}) is inherited from a reconstructed \emph{*phāṭṭakka}, \textquotedblleft{}gate, door\textquotedblright{}; CDIAL lists Punjabi \emph{phāṭak}, \textquotedblleft{}gate\textquotedblright{}, under it (CDIAL 9069).
\end{cousinweb}

\begin{grammarlens}[title={What you've built}]
One more word for this chapter.
\end{grammarlens}

\subsection*{Your turn}
\begin{itemize}
\raggedright
  \item \emph{Say it:} \emph{phāṭak}
  \item \emph{Say it:} \emph{phāṭak}, once more
  \item \emph{Say it:} \emph{jadoṁ rel āundī hai tāṁ phāṭak band hundā hai.} — When the train comes, the crossing gate is closed.
  \item \emph{Recall:} say the Punjabi for across, then the Punjabi for to pass through, then say \emph{phāṭak} again
  \item \emph{Write it:} copy \textbf{\pa{ਫਾਟਕ}} once, checking the whole word
\end{itemize}

\begin{tcolorbox}[breakable,colback=teal!4,colframe=teal!35!black,title={Before you move on}]
What does \pa{ਫਾਟਕ} mean? (\textquotedblleft{}A big gate; a level-crossing gate\textquotedblright{}.) Say it once more.
\end{tcolorbox}
\section[\texorpdfstring{\pa{ਖੰਭਾ} (khambhā)}{ਖੰਭਾ (khambhā)}]{\pa{ਖੰਭਾ} (khambhā) — a pole; a post}

\label{lesson:PA-C264-khambha}



\begin{quote}
Before the new one: say the Punjabi for to pass through, then the Punjabi for a big gate.
\end{quote}

\subsection*{You'll want to know: \pa{ਖੰਭਾ}}
\textbf{\pa{ਖੰਭਾ}} — \emph{khambhā} — \textquotedblleft{}a pole; a post\textquotedblright{}.

\begin{quote}
\pa{ਬੱਸ} \pa{ਬਿਜਲੀ} \pa{ਦੇ} \pa{ਖੰਭੇ} \pa{ਕੋਲ} \pa{ਰੁਕਦੀ} \pa{ਹੈ।}
\end{quote}

\begin{quote}
\emph{bass bijlī de khambhe kol rukdī hai.}
\end{quote}

\begin{quote}
The bus stops by the electricity pole.
\end{quote}

A masculine noun.

\begin{cousinweb}
\textbf{\pa{ਖੰਭਾ}} (\emph{khambhā}) is inherited from Sanskrit \emph{skambha}, \textquotedblleft{}prop, pillar\textquotedblright{}; CDIAL lists Punjabi \emph{khambhā}, \textquotedblleft{}wooden prop, post\textquotedblright{}, under it (CDIAL 13639).
\end{cousinweb}

\begin{grammarlens}[title={What you've built}]
One more word for this chapter.
\end{grammarlens}

\subsection*{Your turn}
\begin{itemize}
\raggedright
  \item \emph{Say it:} \emph{khambhā}
  \item \emph{Say it:} \emph{khambhā}, once more
  \item \emph{Say it:} \emph{bass bijlī de khambhe kol rukdī hai.} — The bus stops by the electricity pole.
  \item \emph{Recall:} say the Punjabi for to pass through, then the Punjabi for a big gate, then say \emph{khambhā} again
  \item \emph{Write it:} copy \textbf{\pa{ਖੰਭਾ}} once, checking the whole word
\end{itemize}

\begin{tcolorbox}[breakable,colback=teal!4,colframe=teal!35!black,title={Before you move on}]
What does \pa{ਖੰਭਾ} mean? (\textquotedblleft{}A pole; a post\textquotedblright{}.) Say it once more.
\end{tcolorbox}
\section[\texorpdfstring{\pa{ਕੋਨਾ} (konā)}{ਕੋਨਾ (konā)}]{\pa{ਕੋਨਾ} (konā) — a corner}

\label{lesson:PA-C264-kona}



\begin{quote}
Before the new one: say the Punjabi for a big gate, then the Punjabi for a pole.
\end{quote}

\subsection*{You'll want to know: \pa{ਕੋਨਾ}}
\textbf{\pa{ਕੋਨਾ}} — \emph{konā} — \textquotedblleft{}a corner\textquotedblright{}.

\begin{quote}
\pa{ਮੇਰਾ} \pa{ਘਰ} \pa{ਗਲੀ} \pa{ਦੇ} \pa{ਕੋਨੇ} \pa{ਉੱਤੇ} \pa{ਹੈ।}
\end{quote}

\begin{quote}
\emph{merā ghar galī de kone utte hai.}
\end{quote}

\begin{quote}
My house is on the corner of the lane.
\end{quote}

A masculine noun.

\begin{cousinweb}
\textbf{\pa{ਕੋਨਾ}} (\emph{konā}) is inherited from Sanskrit \emph{koṇa}, \textquotedblleft{}corner, angle\textquotedblright{}; CDIAL lists Punjabi \emph{koṇā}, \textquotedblleft{}corner\textquotedblright{}, under it (CDIAL 3504).
\end{cousinweb}

\begin{grammarlens}[title={What you've built}]
One more word for this chapter.
\end{grammarlens}

\subsection*{Your turn}
\begin{itemize}
\raggedright
  \item \emph{Say it:} \emph{konā}
  \item \emph{Say it:} \emph{konā}, once more
  \item \emph{Say it:} \emph{merā ghar galī de kone utte hai.} — My house is on the corner of the lane.
  \item \emph{Recall:} say the Punjabi for a big gate, then the Punjabi for a pole, then say \emph{konā} again
  \item \emph{Write it:} copy \textbf{\pa{ਕੋਨਾ}} once, checking the whole word
\end{itemize}

\begin{tcolorbox}[breakable,colback=teal!4,colframe=teal!35!black,title={Before you move on}]
What does \pa{ਕੋਨਾ} mean? (\textquotedblleft{}A corner\textquotedblright{}.) Say it once more.
\end{tcolorbox}
